\subsection{欧拉–庞加莱示性数}

在本节中，我们考虑一个由A-模的类组成的集合C，它是可加的且左正合的，即满足以下两个条件：

\begin{enumerate}
\def\labelenumi{(\Alph{enumi})}
\item
  若M和N是两个属于类型C的A-模，则M ⊕ N属于类型C。
\item
  若 \(0 \to \mathbf{M}' \to \mathbf{M} \to \mathbf{M}'' \to 0\) 是 A-模的一个正合列，且若 \(\mathbf{M}\) 与 \(\mathbf{M}''\) 属于类型 \(\mathcal{C}\) ，则 \(\mathbf{M}'\) 属于类型 \(\mathcal{C}\) .
\end{enumerate}

我们说 \(\mathcal{C}\) 是稳定的，如果它满足以下条件，这些条件蕴含(A)和(G)：

\begin{enumerate}
\def\labelenumi{(\Alph{enumi})}
\setcounter{enumi}{4}
\item
  （“C 在扩张下稳定。”）若 0 → M' → M → M'' → 0 是 A-模的短正合列，且 M' 与 M'' 属于 C 型，则 M 也属于 C 型。
\item
  （“C 在核与余核下稳定。”）对于类型 C 的 A-模的每个同态 f，A-模 Ker f 和 Coker f 均为类型 C。
\end{enumerate}

我们用 K( \(\mathcal{C}\) ) 表示 \(\mathcal{C}\) 的格罗滕迪克群，用 {[}M{]} \(_{\mathcal{C}}\) 或 {[}M{]} 表示由 A-模 M 定义的 K( \(\mathcal{C}\) ) 的元素（VIII，§ 6，第 2 节）。设 G 是一个交换群， \(\varphi\) 是从 K( \(\mathcal{C}\) ) 到 G 中的一个同态。

例子. --- 1) 若 A 是一个域，可取 C 为有限维向量空间的类组成的集合，并取 φ 为 K(C) 到 Z 上的同构，由 φ({[}M{]}) = dim (M) 定义。

\begin{enumerate}
\def\labelenumi{\arabic{enumi})}
\setcounter{enumi}{1}
\tightlist
\item
  可取 \(\mathcal{C}\) 为有限长度模的类构成的集合，并取 \(\varphi: \mathbf{K}(\mathcal{C}) \to \mathbf{Z}\) 为由 \(\varphi([\mathbf{M}]) = \text{long}_{\mathbf{A}}(\mathbf{M})\) 定义的同态。
\end{enumerate}

称分次A-模M为C型的，如果 \(M_{n}\) 对每个 n 而言都是 C 型的（为此，若 M 有界，则这是必要的；若 C 是稳定的，则这是充分的，即模 M 为 C 型）。

定义 8. --- 设 M 为一个 C 型的有界分次 A-模，并设 (M \(_{n}\) ) 为其分次。称 ∑ (-1) \(^{n}\) φ({[}M \(_{n}\) {]}) （G 的元素）为 M 的 φ-示性数，记作 χφ(M) 或简记为 χ(M)。

该定义特别适用于 M 是 A-模复形的底层分次模的情形。

例子. --- 3) 若 M 是 C 型有界的，则对每个 \(p \in Z\) ，M(p) 亦然，并且我们有 \(\chi(\mathbf{M}(p)) = (-1)^{p} \chi(\mathbf{M})\) .

\begin{enumerate}
\def\labelenumi{\arabic{enumi})}
\setcounter{enumi}{3}
\tightlist
\item
  设 \(0 \to \mathbf{M}' \to \mathbf{M} \to \mathbf{M}'' \to 0\) 是分次 A-模与次数为 0 的分次同态的一个正合列。若 \(\mathbf{M}, \mathbf{M}'\) 与 \(\mathbf{M}''\) 是 \(\mathcal{C}\) 型有界的，则我们有
\end{enumerate}

\[
\chi (\mathbf {M}) = \chi (\mathbf {M} ^ {\prime}) + \chi (\mathbf {M} ^ {\prime \prime}).
\]

如果 M 和 M'' 是类型 C 有界的，则 M' 亦然；如果 C 是稳定的，且三个模中有两个是类型 C 有界的，则第三个亦然。

\begin{enumerate}
\def\labelenumi{\arabic{enumi})}
\setcounter{enumi}{4}
\tightlist
\item
  设 \(u: \mathbf{C}' \to \mathbf{C}\) 是 \(\mathcal{C}\) 型有界复形的一个态射。则 Con \((u)\) 是 \(\mathcal{C}\) 型有界的，并且我们有：
\end{enumerate}

\[
\chi (\operatorname{Con} (u)) = \chi (\mathbf {C}) - \chi (\mathbf {C} ^ {\prime}).
\]

\begin{enumerate}
\def\labelenumi{\arabic{enumi})}
\setcounter{enumi}{5}
\tightlist
\item
  可将 G 取为群 K(C) 本身，并取 φ 为恒等映射；此时我们用 \(\chi_{\mathcal{C}}(\mathrm{M})\) 表示 K(C) 的元素 \(\chi_{\varphi}(\mathrm{M}) = \sum (-1)^{n}[\mathrm{M}_{n}]\) 。
\end{enumerate}

注记。--- 称相对于 \(\varphi\) 的元素 \(\mathrm{P}_{\mathbf{M}}(t) = \sum \varphi([\mathrm{M}_{n}]) t^{n} \in \mathrm{G} \otimes \mathbf{Z}[t, t^{-1}]\) 为 M 的庞加莱多项式。我们有 \(\mathrm{P}_{\mathbf{M}}(1) = \varphi([\mathrm{M}])\) 与 \(\mathrm{P}_{\mathbf{M}}(-1) = \chi(\mathrm{M})\) .

引理4. --- 设C是一个有界复形，类型为 \(\mathcal{C}\) 。如果 H(C) = 0，我们有 \(\chi(\mathbf{C}) = 0\) 这由第八卷，第6节，第1号，命题1的推论得出。

命题10. --- 设 C 和 C' 是两个 C 型有界复形。若存在一个拟同构 \(u: \mathbf{C}' \to \mathbf{C}\) ，则我们有 \(\chi(\mathbf{C}) = \chi(\mathbf{C}')\) .

确实，Con \((u)\) 是 \(\mathcal{C}\) 型有界的，并且我们有 \(\chi(\text{Con } (u)) = \chi(\text{C}) - \chi(\text{C}')\) ；另一方面，由 X，第 38 页，推论，H(Con \((u)) = 0\) ，因此 \(\chi(\text{Con } (u)) = 0\) （引理4）。

命题11. --- 设C是一个有界C型复形。

\begin{enumerate}
\def\labelenumi{\alph{enumi})}
\item
  若 \(\mathcal{C}\) 是稳定的，则 H(C) 是 \(\mathcal{C}\) 型的 .
\item
  如果 H(C) 是类型 C，那么 B(C) 和 Z(C) 也是，并且我们有 \(\chi(\mathrm{H}(\mathrm{C})) = \chi(\mathrm{C})\) .
\item
  若 \(\mathcal{C}\) 是稳定的，则对每个 \(n\) ，模 \(\mathbf{Z}_n(\mathbf{C})\) 是 \(\mathcal{C}\) 型的，作为 \(d_n: \mathbf{C}_n \to \mathbf{C}_{n-1}\) 的核；且 \(\mathrm{H}_n(\mathbf{C})\) 是 \(\mathcal{C}\) 型的，作为 \(\mathbf{C}_{n+1} \to \mathbf{Z}_n\) 的余核。另一方面， \(\mathrm{H}_n(\mathbf{C}) = 0\) （只要 \(\mathbf{C}_n = 0\) ）.
\item
  假设 H(C) 是 C 型的。典范正合列：
\end{enumerate}

\[
0 \to Z _ {n} (\mathbf {C}) \to C _ {n} \to B _ {n - 1} (\mathbf {C}) \to 0
\]

\[
0 \to \mathrm{B} _ {n} (\mathrm{C}) \to \mathrm{Z} _ {n} (\mathrm{C}) \to \mathrm{H} _ {n} (\mathrm{C}) \to 0
\]

从 C 的右界出发对 n 作归纳可知， \(Z_{n}(C)\) 与 \(B_{n}(C)\) 对每个 n 都是 C 型的。于是我们有

\[
\chi (\mathrm{C}) = \chi (\mathrm{Z} (\mathrm{C})) + \chi (\mathrm{B} (\mathrm{C}) (- 1)) = \chi (\mathrm{Z} (\mathrm{C})) - \chi (\mathrm{B} (\mathrm{C})) = \chi (\mathrm{H} (\mathrm{C})).
\]

推论。 --- 若 C 是稳定的且 C 是 C 型有界的，则分次模 H(C) 是 C 型有界的，并且我们有 \(\chi(\mathrm{H}(\mathrm{C})) = \chi(\mathrm{C})\) .

命题12. --- 设 \(0 \to C' \to C \to C'' \to 0\) 是复形的一个正合列。a) 若 H(C)、H(C') 和 H(C'') 是 C 型有界的，则我们有

\[
\chi (\mathrm{H} (\mathbf {C})) = \chi (\mathrm{H} (\mathbf {C} ^ {\prime})) + \chi (\mathrm{H} (\mathbf {C} ^ {\prime \prime})) .
\]

\begin{enumerate}
\def\labelenumi{\alph{enumi})}
\setcounter{enumi}{1}
\tightlist
\item
  如果C是稳定的，并且分次模H(C)、H(C')和H(C'')中有两个是C型有界的，那么第三个也是。
\end{enumerate}

a) 由将引理4应用于由所给正合列相关联的同调正合列所定义的零同调复形得出。b) 通过考虑该同调正合列，由以下引理得出：

引理5. --- 设 M → N → P → Q → R 是 A-模的正合列。若 C 是稳定的，且 M、N、Q 和 R 属于 C 型，则模 P 属于 C 型。

令 \(\mathbf{N}' = \operatorname{Coker} (\mathbf{M} \to \mathbf{N})\) 与 \(\mathbf{Q}' = \operatorname{Ker} (\mathbf{Q} \to \mathbf{R})\) 。模 \(\mathbf{N}'\) 与 \(\mathbf{Q}'\) 属于类型 \(\mathcal{C}\) ，并且我们有一个正合列 \(0 \to \mathbf{N}' \to \mathbf{P} \to \mathbf{Q}' \to 0\) .

推论。--- 假设 C 是稳定的，并设 \(u : C' \to C\) 是复形的一个态射，使得 \(\mathrm{H}(\mathbf{C})\) 与 \(\mathrm{H}(\mathbf{C}')\) 是 C 型有界的。则 \(\mathrm{H}(\mathrm{Con}(u))\) 是 C 型有界的，并且我们有

\[
\chi (\mathrm{H} (\text { Con } (u))) = \chi (\mathrm{H} (\mathrm{C})) - \chi (\mathrm{H} (\mathrm{C} ^ {\prime})) .
\]

这由命题 12 应用于复形的正合列（X，第 37 页，命题 7）得出

\[
0 \rightarrow \mathrm{C} \rightarrow \operatorname{Con} (u) \rightarrow \mathrm{C} ^ {\prime} (- 1) \rightarrow 0.
\]

注记。--- 设 E 是一个复形， \(h: \mathrm{E} \to \mathrm{C}\) 与 \(h': \mathrm{E} \to \mathrm{C}'\) 是同伦等价，其中 C 与 C' 是 \(\mathcal{C}\) 型有界的。则我们有 \(\chi(\mathrm{C}) = \chi(\mathrm{C}')\) 。事实上，若 \(h_1\) 是 \(h\) 在同伦意义下的逆，则 \(h' \circ h_1\) 是一个同伦等价，因此是从 C 到

C' 的一个拟同构，从而可以应用命题10。因此，只要存在从 E 到 C 型有界复形 C 的同伦等价，就可以通过设定 \(\chi(E) = \chi(C)\) 来扩展定义 8。命题 10、11、12 及其推论均可推广到这一情形。

应用：

* 设 X 是一个容许有限胞腔分解的拓扑空间（参见第 3 节）。a) 设 K 和 K' 是两个域，令 \(b_i = \dim_K(\mathrm{H}_i(\mathrm{X}, \mathrm{K}))\) 与 \(b_i' = \dim_K(\mathrm{H}_i(\mathrm{X}, \mathrm{K}'))\) 。我们不一定有 \(b_i = b_i'\) ，但有 \(\Sigma(-1)^i b_i = \Sigma(-1)^i b_i'\) .

\begin{enumerate}
\def\labelenumi{\alph{enumi})}
\setcounter{enumi}{1}
\tightlist
\item
  设 \((\mathbf{X}_{n})\) 与 \((\mathbf{X}_{n}^{\prime})\) 是 X 的两个有限胞腔分解，并设 \(c_{n}\) 与 \(c_{n}^{\prime}\) 是这两个分解中 n 维胞腔的个数。我们有
\end{enumerate}

\[
\Sigma (- 1) ^ {i} c _ {i} = \Sigma (- 1) ^ {i} c _ {i} ^ {\prime}.
\]

\begin{enumerate}
\def\labelenumi{\alph{enumi})}
\setcounter{enumi}{2}
\tightlist
\item
  利用 \(a)\) 与 \(b)\) 的记号，我们有 \(\Sigma(-1)^{i} c_{i} = \Sigma(-1)^{i} b_{i}\) .
\end{enumerate}

性质 \(a\) ) 与 \(b\) ) 由 \(c\) ) 推出，而 \(c\) ) 由命题11应用于第 3 号中所述的复形 \(\Gamma\) 得出，取 \(\mathscr{C}\) 为有限维 K-向量空间的类，取 \(\varphi\) 为由 \(\varphi([\mathrm{M}]) = \dim_{\mathrm{K}}(\mathrm{M})\) 定义的函数（X，第 40 页，例 1）。*

